\documentclass[dvipdfmx]{jlreq}

\title{アドカレ}
\date{\today}

\begin{document}
\maketitle

直腸圧$p_s$は、粘性無しでも要する圧力$p_a$と流体の粘性に抗する圧力$p_b$の和によって記述されます。
\[
p_s = p_a + p_b
\]
ここで、糞は完全流体と近似され、$p_a$は糞を初速$v$まで加速するために必要な撃力として、以下のように糞の密度$\rho$と$v$を用いて見積もられました。
\[
p_a= \rho v^2
\]
また、ハーゲン・ポアズイユ流(円管内の定常層流)を仮定した解析から、粘性補正項$p_b$は糞の体積$V$, 長さ$l$, 粘度$\eta$, 地面に落ちるまでの時間$t$, 排泄腔の半径$r$を用いて以下のように評価されました。
\[
p_b= \frac{8Vl\eta}{\pi r^4 t}
\]

力学の斜方投射の問題と同様に、初期条件$(x_0,y_0)=(0,h), (v_{0x},v_{0y})=(v_0 \cos\theta, v_0 \sin\theta)$のもとで運動方程式
\[
m \frac{d^2x}{dt^2} = 0
\]
\[
m \frac{d^2y}{dt^2} = -mg
\]
を解くと、糞の到達する長さ$d=v_0 t \cos\theta$は
\[
d=v_0 \cos\theta \bigg(
\frac{v_0\sin \theta}{g}
+ \sqrt{\frac{v^2_0\sin^2 \theta}{g^2}+\frac{2h}{g}}
\bigg)
\]
と求まり、その最大値は
\[
d_{max}=v_0 \cos \Bigg[\sin^{-1}\bigg( 
\frac{1}{\sqrt{2\Big(1+\frac{hg}{v^2_0} \Big)}}
\bigg) \Bigg]
\Bigg(
\frac{v_0}{g} \frac{1}{\sqrt{2\Big(1+\frac{hg}{v^2_0}\Big)}}
+ \sqrt{\frac{v^2_0}{2g^2\Big(1+\frac{hg}{v^2_0}\Big)}+\frac{2h}{g}}
\Bigg)
\]
と評価されました。

Meyer-RochowとGalの論文では、$p_a$は初速度に対応する撃力として導入されましたが、田島と藤澤の論文ではこれをベルヌーイの法則に基づき評価し直しました。胃を半径$R$の円筒とみなし、モノがすべて腹の中にあるときの液面の高さを$z$とおくと、圧力$P$について
\[
\frac{1}{2}\rho v^2 + \rho gz + P = const.
\]
というベルヌーイの法則が成り立つので、発射の前後で
\[
\rho gz + P_a =\frac{1}{2}\rho v^2_0
\]
という関係が成り立ちます(大文字の$P_a$と書いて先行研究の表記と区別しました)。排出した液量$V=\pi z R^2$から
\[
P_a = \frac{1}{2}\rho v^2_0 - \frac{\rho g}{\pi R^2} V
\]
と$P_a$が求まります。$V$は軌道の長さ$l$(先行研究では糞の届いた距離で近似していましたが、ここでは正確に取り扱います)と噴出後の流れの断面積$S$の積で書けますが、一般に円筒状の$S=\pi r^2$とはならず、
\[
S= C \pi r^2
\]
と収縮率$C$を用いて書けることが知られています。

加えて、粘性補正の寄与についても見直しが図られました。ハーゲン・ポアズイユ流を仮定して、地面に落ちるまでの時間$t$から計算した寄与$P_b$は変わりませんが、胃腸の中にある時間$t_d \approx \frac{z}{(0+v_0)/2}$から計算できる寄与
\[
P_c =\frac{8Vz\eta}{\pi r^4 t_d} \approx \frac{4Vv_0\eta}{\pi r^4}
\]
も含めた$P=P_a+P_b+P_c$が計算されました。

固液混合物の溶媒に対する相対粘度について、次のような経験式(クリーガーの式)
\[
\eta_r = (1-\frac{\phi}{\phi_{m}})^{-k\phi_{m}} 
\]
が知られています。ここで、$\phi_{(m)}$は(最大)体積分率、$k$は実験定数です。
様々な排泄物の組成からこの式を逆に解くことで、パラメータ$k,\phi_{m}$を実験的に求められるかもしれません。

また、固体粒子の含まれる割合の高さを加味すると、排泄物は一般に非ニュートン液体として扱うべきである可能性があります。たとえば、せん断応力が閾値(降伏応力)を超えるまで流動しない塑性流体は、ペンギンの排泄物のモデルとして適切かもしれません。

このようなフン体のレオロジーは、物理学に残されたフロンティアの一つだと言えるでしょう。


\end{document}
